%%%%%%%%%%%%%%%%%%%%%%%%%%%%%%%%%%%%%%%%%%%%%%%%%%%%%%%%%%%%%%%%%%%%%%%%%%%%%%%
\chapter{Introduction} %%%%%%%%%%%%%%%%%%%%%%%%%%%%%%%%%%%%%%%%%%%%%%%%%%%%%%%%

\section{Motivation}

Rowing is a technically demanding sport. Good performance does not only depend on strength or endurance, but also on how well the athlete coordinates the movement during each stroke. A rowing stroke is usually described by the phases catch, drive, finish, and recovery. In these phases the athlete, the sliding seat, the oars, and the boat form one coupled mechanical system. Small changes in timing, force application, body position, or boat balance can influence the efficiency of the stroke. Reviews of rowing technique show that coaches and researchers have long tried to connect stroke mechanics with boat velocity and performance, but also that it is difficult to define one single ideal technique that works for every athlete and setup \cite{soper_towards_2004}.

This makes rowing a good use case for objective sensing. A coach can see many errors, but visual feedback is still subjective and depends on experience. Video can help, but it needs time, good camera placement, and manual interpretation. In many normal training situations, especially outside elite clubs, athletes do not have constant access to detailed feedback. A small sensing system that can collect repeatable movement data could therefore support training and later analysis.

At the same time, such a system must not disturb the athlete. Rowing is sensitive to small changes in rhythm and posture. A measurement device that is heavy, visible, wired, or difficult to mount can change the situation it is supposed to measure. Llosa et al. make this point clearly when discussing wired rowing measurement equipment: if athletes know they are being measured because the system is visible and intrusive, their behavior can change and distort the result \cite{llosa_remote_2009}. This is an important design requirement for this thesis. The prototype should be small, low-cost, and unobtrusive enough that it can be used in realistic training without turning the training session into a laboratory experiment.

Inertial measurement units (IMUs) are attractive for this problem because they are small, low-cost, and technically mature. They can record acceleration and angular velocity directly on the athlete, the boat, or the equipment. IMUs are already used in many sport and movement-analysis systems because they do not require fixed cameras or a prepared laboratory space \cite{arlotti_benefits_2022}. In rowing, IMUs have been used to measure boat motion, oar motion, handle motion, athlete movement, and rowing-related feedback variables \cite{llosa_remote_2009,junker_sensor-based_2015,hohmuth_wireless_2023,rodriguez_mendoza_wearable_2024}. This makes them a strong candidate for a low-cost rowing prototype that can be used outside the lab.

However, IMUs also have a central weakness. They do not directly measure position. Position must be estimated by integrating acceleration, and small sensor errors, bias, orientation errors, and noise grow over time during this integration. In inertial navigation, this drift is a basic problem. Rodriguez Mendoza and O'Keefe explicitly describe that inertial navigation solutions degrade over time because accelerations and angular rates are integrated to determine velocity, position, and attitude \cite{rodriguez_mendoza_wearable_2024}. For this reason, many advanced systems use sensor fusion, for example by combining inertial sensors with GNSS, UWB, external reference systems, or biomechanical models \cite{cloud_adaptive_2019,rodriguez_mendoza_wearable_2024,ruffaldi_sensor_2015}.

This thesis uses a simpler and more robust idea. Instead of trying to reconstruct the exact absolute position of the rowing seat, the prototype focuses on relative and stroke-level metrics. The main target is not to say: ``the seat is exactly at this position in meters''. The target is to measure signals that are useful for rowing feedback, such as forward acceleration, stroke timing, movement intensity, pitch and yaw stability, and the difference between a seat-mounted sensor and a boat-mounted reference sensor. In the two-sensor case, the core idea is:

\begin{equation}
  \text{relative motion} = \text{SEAT signal} - \text{BOAT signal}.
\end{equation}

This is important because a seat-mounted IMU on its own measures both the athlete's motion and the motion of the boat under the athlete. A boat-mounted IMU can provide a reference for the boat motion. Subtracting the boat signal from the seat signal can therefore help isolate movement that is more related to the athlete and sliding seat. This principle is close to the seat-position approach by Gravenhorst et al., who used one sensor on the sliding seat and one sensor on the boat shell to calculate seat acceleration relative to the boat before further processing \cite{gravenhorst_self-aligning_2012}. The difference in this thesis is that the prototype does not need to solve the full problem of accurate double-integrated seat position. It can avoid part of the drift problem by using relative metrics and per-stroke signal features.

\section{Problem Context}

The problem addressed in this thesis is not only a software problem and not only a hardware problem. It sits between embedded sensing, sport-specific requirements, signal processing, and physical integration. A useful rowing sensor has to record data at a sufficient sampling rate, store or transmit it reliably, survive the rowing environment, and be mounted in a repeatable way. If the sensor is loose, badly sealed, or placed differently in every session, the measured signal can change even if the athlete's technique stays the same.

The rowing environment is harder than many indoor movement-analysis tasks. The boat is not a fixed reference frame. It accelerates and decelerates during the stroke. It also moves due to water, wind, steering, and the athlete's body mass. A sensor on the athlete or seat therefore observes a mixture of several motions. If the goal were exact absolute seat position, the system would need careful calibration, synchronization, drift correction, and validation against a reference system. For a bachelor-level prototype, this would be a high-risk target and would require more time than the development of the embedded prototype itself.

For this reason, the thesis uses a more modest but still useful goal. The prototype is designed to collect reliable IMU data and derive relative metrics that can support rowing feedback. These metrics are easier to obtain than absolute position because they can be based on signal shape, peaks, timing, standard deviation, root mean square values, and relative differences between sensors. Such metrics still describe meaningful parts of the rowing motion, especially when they are compared between strokes, between sessions, or between a seat sensor and a boat reference sensor.

This also changes how the evaluation should be understood. The system is not evaluated as a precise motion-capture replacement. It is evaluated as a prototype that can record useful signals, produce interpretable metrics, and form a basis for later validation. This is a realistic scope for a low-cost wearable sensing approach.

\section{Aim of the Thesis}

The aim of this thesis is to design and evaluate a low-cost IMU-based wearable sensing prototype for rowing. The prototype should capture movement data that is relevant to rowing technique analysis while staying simple enough to be built, mounted, and tested in a realistic environment.

The thesis focuses on four main goals:

\begin{itemize}
  \item to define rowing-specific requirements for a small IMU-based sensing system;
  \item to implement an embedded prototype that records and streams rowing-related IMU data;
  \item to develop an analysis pipeline for stroke-level and relative metrics;
  \item to consider hardware integration and enclosure design as part of data quality.
\end{itemize}

The prototype is based on a small embedded sensing unit with an IMU, SD card logging, USB serial output, and Bluetooth Low Energy (BLE) streaming. The current embedded software records acceleration and angular velocity data at approximately 100 Hz. The data are written as comma-separated values and include a device identifier, sequence number, timestamp, acceleration axes, and gyroscope axes. The software supports a device role such as \texttt{SEAT}, and the same code can be adapted for a second unit such as \texttt{BOAT}.

The software side contains three main tools. First, a live serial or BLE reader can receive the CSV stream and calculate simple rolling metrics. Second, a live web dashboard can display the current signal, record a training segment, and mark good or bad strokes. Third, an offline analysis script can load a recorded CSV file, check the sample timing, detect sequence gaps, calculate summary statistics, estimate stroke timing, and generate an HTML report. In the two-IMU mode, the live dashboard can connect to a seat unit and a boat unit and compute relative signals from the difference between both streams.

\section{Research Questions}

The following research questions guide the work:

\begin{enumerate}
  \item Which requirements should a low-cost IMU-based rowing sensing prototype fulfill with respect to sensing, mounting, robustness, and usability?
  \item How can a small embedded system record rowing-related acceleration and angular velocity data in a practical and repeatable format?
  \item Which stroke-level metrics can be derived from IMU signals without claiming exact absolute seat position?
  \item How can a second boat-mounted IMU be used as a reference to compute relative seat-to-boat movement indicators?
  \item Which limitations remain because of IMU drift, mounting quality, synchronization, and lack of high-precision ground truth?
\end{enumerate}

These questions are intentionally practical. The goal is not to build the most complex rowing measurement system. The goal is to build a useful and honest prototype that can be improved in later work.

\section{Design Philosophy}

The main design choice is to treat the IMU data as a source for robust training indicators rather than as a perfect position sensor. The analysis therefore uses the mounted X axis as a forward/backward candidate axis, smooths and centers the signal for stroke detection, and calculates stroke rate and movement proxies. A relative velocity proxy can be estimated by integrating a smoothed acceleration segment with simple drift correction, but it is clearly treated as a proxy and not as an absolute speed in meters per second. A power-transfer proxy can be calculated from the product of forward acceleration and relative speed, but it is also not interpreted as real mechanical power in watts.

This approach is intentionally conservative. It accepts the known limits of low-cost IMUs and avoids overclaiming. At the same time, it keeps the system useful for a first prototype. The system can still show whether the movement is periodic, whether strokes are detected consistently, whether the boat reference signal differs from the seat signal, and whether relative movement patterns are stable across strokes.

This design philosophy also makes the prototype more affordable than many alternatives. Systems that use UWB, GNSS, force sensors, EMG, optical motion capture, or full-body IMU networks can measure more variables, but they are more expensive and more complex to install \cite{hohmuth_wireless_2023,rodriguez_mendoza_wearable_2024,ruffaldi_sensor_2015}. The proposed prototype sacrifices some measurement richness in exchange for lower cost, lower setup effort, and simpler interpretation.

\section{Contributions}

The thesis contributes a practical prototype and analysis pipeline for low-cost rowing motion sensing. The main contributions are:

\begin{itemize}
  \item an embedded IMU logger for rowing-related acceleration and angular velocity data;
  \item a simple CSV-based data format with timestamps and sequence numbers for later analysis;
  \item USB serial and BLE streaming for live monitoring;
  \item an offline analysis workflow for sampling checks, signal statistics, stroke detection, and report generation;
  \item a two-IMU relative measurement concept using a seat-mounted and boat-mounted sensor;
  \item a metric strategy that focuses on relative and stroke-level indicators instead of exact absolute seat position;
  \item a hardware-integration perspective that treats enclosure design, mounting, and robustness as part of measurement quality.
\end{itemize}

The expected result is not a final commercial rowing coach. It is a research prototype that can be evaluated, improved, and used to collect first data. The value of the work lies in connecting requirements, hardware design, embedded sensing, data processing, and rowing-specific interpretation in one system.

\section{Thesis Structure}

The rest of this thesis is structured as follows. Chapter~2 reviews related work on rowing technique, IMU-based rowing systems, seat tracking, sensor fusion, general IMU limitations, and enclosure design. Chapter~3 describes the system requirements and design decisions. Chapter~4 explains the implementation of the embedded logger, data pipeline, live dashboard, and offline analysis. Chapter~5 presents the prototype case and experimental evaluation. Chapter~6 discusses the results, limitations, and future work. Chapter~7 concludes the thesis.


%%%%%%%%%%%%%%%%%%%%%%%%%%%%%%%%%%%%%%%%%%%%%%%%%%%%%%%%%%%%%%%%%%%%%%%%%%%%%%%
\chapter{Related Work} %%%%%%%%%%%%%%%%%%%%%%%%%%%%%%%%%%%%%%%%%%%%%%%%%%%%%%%%

\section{Overview}

This chapter reviews related work from rowing measurement systems, IMU-based motion analysis, wearable sensing, and waterproof prototype design. The goal is not only to list existing systems. The goal is to understand which design decisions are useful for this thesis and which approaches are too complex, too expensive, too intrusive, or too dependent on exact position estimation.

The reviewed work shows a clear pattern. Rowing systems can measure many different quantities: boat speed, boat acceleration, oar angle, handle force, pressure distribution, seat displacement, muscle activity, synchrony, and 3D handle motion \cite{llosa_remote_2009,avvenuti_mars_2013,cloud_adaptive_2019,hohmuth_wireless_2023,rodriguez_mendoza_wearable_2024}. Rich systems can provide more information, but they also add cost, setup time, calibration effort, and possible disturbance to the athlete. In contrast, simple systems are easier to use but usually measure fewer variables. This thesis is positioned on the simple side of this trade-off: it aims to provide useful relative metrics with low-cost IMUs instead of full biomechanical reconstruction.

\section{Rowing Technique and the Need for Objective Feedback}

Rowing performance depends on a coordinated stroke cycle. In sculling and sweep rowing, the athlete must move the body, seat, oars, and boat in a controlled sequence. Soper and Hume describe rowing performance as a complex interaction of kinematics, kinetics, anthropometry, boat setup, and coaching practice \cite{soper_towards_2004}. They also point out that rowing biomechanics has investigated many aspects of technique, but that simple and general predictors of performance are difficult to define. This is important for the present thesis because it shows that a sensing system should not claim to measure ``perfect technique'' directly. It should instead provide objective data that can help interpret technique.

Hohmuth et al. also describe rowing as a sport in which technique, timing, force, and coordination strongly influence performance \cite{hohmuth_wireless_2023}. Their work lists typical technique-related variables such as stroke length, blade depth, catch and finish timing, smooth recovery, force application, seat displacement, and boat acceleration. This supports the idea that a useful rowing sensing prototype should not only measure one single number. It should allow several views of the stroke, such as timing, forward movement, stability, and periodic signal shape.

Feedback systems in rowing are especially valuable because the athlete cannot easily see the own movement while rowing. Coaches often use observation and video, but objective sensor data can help make repeated training sessions more comparable. Junker's doctoral work on sensor-based performance feedback in rowing shows the broader research direction of using sensors to create quantitative feedback for rowing technique \cite{junker_sensor-based_2015}. This thesis follows the same general direction, but with a smaller and cheaper prototype focused on IMU-based relative metrics.

\section{Observation, Video, and Commercial Instruments}

The simplest approach to rowing feedback is direct observation by a coach. This has a major advantage: the coach can see the full athlete and can interpret the movement in context. However, observation is subjective and hard to compare across sessions. Video can make the feedback more repeatable, but it still requires camera setup, synchronization, manual review, and expertise. Video also does not directly provide acceleration or angular velocity signals.

Commercial rowing instruments can measure useful variables such as boat speed, stroke rate, force, and oar angle. They can be very valuable for training. However, they are often designed for specific equipment, can be expensive, and can require installation effort. Llosa et al. compare their wireless REMOTE system with RowX, a commercial wired instrument, and describe that wired infrastructure can be bulky and difficult to install and uninstall \cite{llosa_remote_2009}. They also state that visible measurement equipment can change athlete behavior because the athletes know they are being measured. This point is important for the present thesis. A measurement device should be as unobtrusive as possible, because a device that changes posture or rhythm can reduce the validity of the measurement.

This motivates a low-cost and compact prototype. A small IMU-based unit cannot measure everything that a commercial system can measure, but it can reduce cost, wiring, and setup effort. It also gives access to raw movement data, which can be processed later.

\section{Wireless Sensor Networks for Rowing}

One early example of a wireless rowing sensing system is REMOTE by Llosa et al. \cite{llosa_remote_2009}. REMOTE uses wireless sensor network nodes with accelerometers to monitor oar and boat motion. The system was tested through laboratory calibration, ergometer tests, and on-boat tests. The work shows that accelerometer-based wireless systems can provide useful information such as boat motion, oar motion, stroke angles, boat balance, and stroke trajectory.

REMOTE is relevant for this thesis for two reasons. First, it shows that low-cost wireless sensors can be useful for rowing. Second, it shows the practical limits of such systems. The authors attribute deviations in stroke trajectory to handmade packaging, unknown original position, and numerical integration errors \cite{llosa_remote_2009}. These are exactly the kinds of problems that appear when low-cost inertial sensors are used in a real environment. The packaging matters, the starting position matters, and integration drift matters.

The present thesis responds to these problems by not making stroke trajectory or exact seat position the main claim. Instead, it focuses on relative and stroke-level metrics that are less dependent on long integration chains. It also treats enclosure and mounting as part of the measurement problem.

\section{Oar-Based IMU Systems}

Several rowing systems place sensors on the oar or handle. This is useful because oar motion is directly connected to stroke length, catch timing, finish timing, and blade handling. Gravenhorst et al. validated an IMU-based rowing oar angle measurement system and compared the IMU result with a reference system based on a potentiometer \cite{gravenhorst_validation_2013}. This shows that IMUs can be used successfully when the measurement target is well defined and the sensor is mounted in a stable way.

Oar-based systems have an important advantage: they measure a central output of the rowing motion. A coach often cares about oar angle and blade movement because these variables connect directly to propulsion. However, oar motion is also the result of many body movements. If an oar signal is wrong, it can be hard to know whether the cause is leg drive, body swing, arm movement, timing, or boat movement. SonicSeat makes a similar point: knowing the resulting oar movement does not always explain which sub-movement caused the deviation \cite{gravenhorst_sonicseat_2012}.

For this thesis, this means that an oar IMU would not be the best first target. The aim is not to reconstruct the oar path. The aim is to capture a signal that is closer to the seat and leg-driven part of the rowing motion. A seat-mounted IMU is therefore a good practical target, especially when paired with a boat-mounted reference sensor.

\section{Seat Motion as a Rowing Metric}

The sliding seat is an important part of rowing technique because it is directly connected to leg drive and body movement. Gravenhorst et al. argue that seat movement is linked to the rower's leg movement and therefore to one of the most powerful parts of the rowing stroke \cite{gravenhorst_self-aligning_2012}. SonicSeat also emphasizes that the sliding seat allows the rower to use the legs for a longer and more powerful stroke and that seat movement is directly linked to the movement of the rower's center of gravity \cite{gravenhorst_sonicseat_2012}.

Seat motion is attractive for sensing because it is periodic and strongly connected to the stroke cycle. It can support metrics such as stroke rate, catch delay, maximum sliding velocity, and sliding seat displacement \cite{gravenhorst_sonicseat_2012}. These metrics are easier for coaches to interpret than raw accelerometer signals.

However, seat motion is also hard to measure with a single IMU. A sensor on the seat does not only measure the seat moving on the rail. It also measures the boat moving under the seat. This is the central reason why the present thesis uses the idea of a boat reference sensor.

\section{Two-Sensor Seat-to-Boat Measurement}

The most important related work for this thesis is the self-aligning and drift-compensated rowing seat position measurement system by Gravenhorst et al. \cite{gravenhorst_self-aligning_2012}. Their system uses two miniaturized sensor units: one attached to the sliding seat and one attached to the boat shell. The seat sensor measures the superposition of boat acceleration and seat acceleration relative to the boat. The boat sensor measures boat acceleration. The relative seat acceleration can therefore be calculated from the difference between both sensors, after correcting for sensor misalignment.

This is the strongest prior support for the design of this thesis. It gives a clear reason why two sensors are better than one for seat-related rowing metrics. A single seat sensor mixes seat motion and boat motion. A boat sensor can act as a reference. The same basic idea is used here:

\begin{equation}
  a_{\text{relative}}(t) = a_{\text{seat}}(t) - a_{\text{boat}}(t).
\end{equation}

The difference is the target output. Gravenhorst et al. continue toward drift-compensated seat position by double integration, using additional discrete correction events to reduce low-frequency drift \cite{gravenhorst_self-aligning_2012}. The present thesis does not attempt to fully reproduce that position-tracking system. It uses the relative idea but focuses on signals and metrics that can be extracted without exact absolute position. This is a deliberate simplification. It makes the system less powerful, but also less fragile.

\section{Ultrasonic Seat Tracking}

Gravenhorst et al. also developed SonicSeat, a contactless seat position tracker based on ultrasonic sound measurements \cite{gravenhorst_sonicseat_2012}. SonicSeat measures seat position directly with an ultrasonic distance sensor mounted in the boat. It was used to derive metrics such as stroke rate, catch delay, maximum sliding velocity, and sliding seat displacement. The system was also connected to video analysis and expert evaluation.

Ultrasonic sensing has a clear advantage over IMU integration: it can measure distance more directly. This avoids some drift problems. For exact seat displacement, this can be better than a simple IMU. However, it also requires a dedicated distance sensor, stable geometry, line of sight to the seat, and careful mounting. It measures the seat, but it does not automatically provide rotation rates, pitch/yaw stability, or general acceleration signals. It is therefore a strong seat-position solution, but less general than an IMU.

For the present thesis, SonicSeat is useful as a metric reference, not as a hardware model. It shows which seat-related features can matter, but the prototype uses IMUs because they are cheap, compact, and can measure multiple axes of motion.

\section{Multi-Sensor Rowing Systems}

Hohmuth et al. developed WiRMS, a wireless rowing measurement system that combines several sensing modalities \cite{hohmuth_wireless_2023}. Their system includes handle units with IMUs, strain gauges, and pressure sensors, a boat kinematic measurement unit with acceleration and GPS, contactless seat displacement measurement, and synchronized EMG data. A proof-of-concept study with experienced athletes showed that the system can record movement and force variables during on-water rowing.

WiRMS shows what is possible with a rich sensor system. It can measure force, pressure distribution, oar angles, seat displacement, boat acceleration, boat velocity, and muscle activity. This is much more complete than the system in this thesis. However, this completeness has a cost. It needs modified sculls, several sensors, a tablet computer, EMG equipment, synchronization, and more complex setup. The authors also report that some participant data had to be excluded due to an IMU configuration error in the sculls \cite{hohmuth_wireless_2023}. This illustrates a general issue: the more complex the system becomes, the more points of failure it has.

The present thesis takes a different path. It does not try to measure every relevant variable. It tries to build a small prototype that can be understood, mounted, and evaluated within the scope of a bachelor thesis. This makes the system less complete than WiRMS, but also more accessible.

\section{Crew Synchrony and Pattern-Based Feedback}

Avvenuti et al. present MARS, a multi-agent system for assessing rower coordination using motion-based stigmergy \cite{avvenuti_mars_2013}. Their system focuses on crew synchrony rather than exact biomechanical reconstruction. It processes motion signals and produces asynchrony indicators for feedback.

This work is important because it shows that rowing feedback does not always require full 3D reconstruction. Useful feedback can also be based on timing, similarity, and repeated signal patterns. This supports the metric strategy in the present thesis. If the prototype can detect strokes, compare signal shapes, and show relative timing or stability indicators, it can still be useful even without exact position.

For a low-cost prototype, pattern-based metrics are often more realistic than full kinematic reconstruction. They can also be easier to validate with video or coach labels. The live dashboard in this thesis supports this direction by allowing good and bad stroke markers during recordings.

\section{Smartphone-Based Boat Motion and Sensor Fusion}

Cloud et al. use smartphone GPS and accelerometer data to estimate boat speed, distance, and distance per stroke \cite{cloud_adaptive_2019}. They compare a complementary filter and a Kalman filter against differential GPS reference data. Their results show that activity-specific sensor fusion can improve estimates from low-cost consumer sensors. They also show that exact stroke-level distance measurement is difficult with low-cost sensors.

This is relevant for two reasons. First, it supports the idea that low-cost sensors can become more useful when the processing is specific to rowing. Second, it warns against overclaiming. Even with sensor fusion, Cloud et al. report that distance-per-stroke accuracy remains limited compared with the needs of very detailed stroke-level analysis \cite{cloud_adaptive_2019}. Therefore, a bachelor prototype based only on low-cost IMUs should be careful with exact distance or position claims.

The present thesis uses this lesson directly. It uses simple signal processing and relative metrics instead of trying to make the IMU act like a high-precision position system.

\section{UWB, GNSS, and High-Accuracy Positioning}

Rodriguez Mendoza and O'Keefe present a wearable multi-sensor positioning prototype for rowing technique evaluation \cite{rodriguez_mendoza_wearable_2024}. Their system combines an IMU, ultra-wideband ranging, and GNSS to estimate the position, velocity, and attitude of a wearable node during rowing. They compare several methods, including UWB, inertial navigation, rowing dead reckoning, GNSS-aided inertial navigation, and UWB-aided inertial navigation.

This approach can provide richer and more accurate position information than a simple IMU-only prototype. However, it requires more hardware and more infrastructure. UWB needs anchors or transmitters. GNSS-based reference solutions require good satellite conditions and, for high accuracy, can require more expensive equipment or post-processing. The system also needs coordinate-frame definitions and synchronization between components.

For this thesis, such a system is too complex as a first prototype. It is valuable as a reference for what high-end rowing positioning can look like, but the design goal here is different. The aim is not to build an outdoor positioning system. The aim is to build a low-cost sensor that can provide relative rowing indicators with minimal setup.

\section{Full-Body and Articulated Tracking}

Ruffaldi et al. investigate sensor fusion for complex articulated body tracking in rowing \cite{ruffaldi_sensor_2015}. Their work represents a more complete body-tracking direction, using several IMUs and constraints to estimate articulated motion. Such a system can provide richer information about upper body and arm movement than a seat/boat prototype.

However, full-body tracking requires several sensors, calibration, a body model, and often a controlled validation environment. It is powerful, but it is not lightweight. For practical rowing training, this can be a barrier. More sensors also mean more mounting positions, more possible placement errors, and more time before training can start.

The present thesis chooses a smaller scope. A seat-mounted and optional boat-mounted IMU cannot explain the full body movement, but it can focus on a repeatable and important part of the stroke. This is a reasonable trade-off for a low-cost wearable prototype.

\section{General Challenges of IMU-Based Motion Analysis}

The main technical challenge of IMU-based motion analysis is drift. Accelerometers and gyroscopes are useful because they provide high-rate motion data, but their errors accumulate when the data are integrated. Bias, noise, temperature effects, wrong initial orientation, and mounting errors can all affect the result. This is why many IMU systems use additional sensors, constraints, or event-based corrections. In rowing, the periodic structure of the stroke can help because peaks or repeated phases can be used for segmentation or reset logic \cite{cloud_adaptive_2019,rodriguez_mendoza_wearable_2024}.

Other sports show the same pattern. Yu et al. use an IMU-based system for football kick motion analysis and show a processing chain including calibration, orientation estimation, and validation against camera-based data \cite{yu_motion_2022}. He et al. validate an IMU-based gait analysis system against optical motion capture, which is a common way to test whether wearable movement metrics are accurate enough \cite{he_accuracy_2024}. These works are not rowing systems, but they show that IMU metrics need careful validation and that accuracy claims should be tied to the measured variable.

Reviews of IMU use in sports and sports medicine also describe the practical benefits and limitations of wearable sensors. Arlotti et al. discuss IMU-based wearables as useful tools for sport and medical contexts, but the broader field still has to handle calibration, drift, placement, and interpretation issues \cite{arlotti_benefits_2022}. Presicci et al. review IMU-based systems in skiing and show that IMUs are widely used for sport-specific metrics, but that protocols and validation still matter strongly \cite{presicci_imu-based_2025}. These lessons transfer to rowing: a metric can be useful only if the mounting, processing, and interpretation are clear.

This thesis therefore uses a cautious interpretation of IMU data. It does not present integrated acceleration as exact seat position. It uses acceleration, angular velocity, signal variation, stroke events, and SEAT--minus--BOAT differences as relative indicators. This reduces the dependence on long-term integration and makes the prototype more suitable for early evaluation.

\section{Wearable Integration and Sensor Placement}

Wearable sensing is not only an algorithmic problem. The sensor must be placed and mounted in a way that is comfortable, repeatable, and robust. Gei\ss{}ler et al. show an alternative wearable sensing approach using textile capacitive sensors embedded in garments \cite{geisler_embedding_2024}. Their work is not an IMU rowing system, but it highlights an important general point: sensing quality is connected to the physical integration of the sensor into a wearable object. If the sensor moves relative to the body or object, the measured signal changes.

This is also true for a rowing IMU prototype. A sensor on the seat or boat must be attached firmly enough that the signal represents the motion of the seat or boat, not vibration of the casing. At the same time, the device should be small and unobtrusive. It must survive water, splashes, mechanical stress, and handling before and after training.

The choice of placement also influences what can be measured. A sensor on the oar is good for oar angle, but not directly for seat motion. A sensor on the boat is good for boat acceleration and reference motion, but not directly for athlete motion. A sensor on the seat is close to the leg-driven movement, but it also contains boat motion. This is why the two-sensor design is attractive: it combines a seat signal and a boat reference while still keeping the system small.

\section{Enclosure Design and Robustness}

Waterproof enclosure design is part of the measurement system. If water enters the enclosure, the electronics can fail. If the enclosure is too large or heavy, it can disturb the athlete or the boat. If the mounting is unstable, the signal quality becomes worse. Therefore, hardware integration is not separate from data quality.

Lindqvist and Stroemberg Groendahl describe a waterproof probe enclosure development process with requirements, concept selection, O-ring sealing, material choices, and prototype testing \cite{lindqvist_design_2015}. Huayna Fernandez similarly focuses on enclosure development for a water-quality measuring tool and discusses the design of a housing for use in a wet environment \cite{huayna_fernandez_enclosure_2025}. These works are not rowing-specific, but they are useful because rowing also exposes electronics to water, handling, vibration, and outdoor conditions.

For this thesis, enclosure design should therefore be treated as a real engineering requirement. A prototype that works on a desk but fails near water is not suitable for rowing. A prototype that measures well only when held by hand is also not suitable. The enclosure and mounting concept must support repeatable sensor orientation and stable attachment.

\section{AI and Future Feedback Systems}

Sorysz discusses opportunities and challenges for artificial intelligence in wearables for physical therapy and sports training \cite{sorysz_artificial_2024}. For this thesis, this does not mean that a complex AI model is required. Instead, it shows that wearable sports systems should be designed with later feedback and interpretation in mind. A clean data format, stable mounting, and meaningful labels such as good or bad strokes can support future learning-based analysis.

This is relevant because the prototype already records structured CSV files and can store markers during live recordings. In later work, these markers could be used to train or evaluate classification models. However, such models would only be meaningful if the underlying data are collected consistently. This again supports the focus on simple hardware, repeatable mounting, and honest metrics.

\section{Comparison of Approaches}

Table~\ref{tab:approach-comparison} summarizes the main approaches discussed in this chapter and compares them to the direction of this thesis.

\begin{table}[htbp]
  \centering
  \caption{Comparison of rowing sensing approaches and their relevance for this thesis.}
  \label{tab:approach-comparison}
  \begin{tabular}{p{0.24\textwidth} p{0.28\textwidth} p{0.34\textwidth}}
    \hline
    Approach & Strength & Limitation for this thesis \\
    \hline
    Coach observation and video & Full-body context, easy to understand & Subjective, time-consuming, not directly numerical \\
    Commercial wired systems & Established rowing metrics & Can be expensive, visible, and intrusive \cite{llosa_remote_2009} \\
    Oar-mounted IMUs & Good for oar angles and stroke length & Does not directly isolate seat or athlete motion \cite{gravenhorst_validation_2013} \\
    Ultrasonic seat tracking & Direct seat distance measurement & Needs dedicated distance hardware and stable geometry \cite{gravenhorst_sonicseat_2012} \\
    Multi-sensor systems & Rich variables such as force, EMG, GPS, seat displacement & More complex, more expensive, more setup effort \cite{hohmuth_wireless_2023} \\
    UWB/GNSS positioning & Higher accuracy position estimates & Needs infrastructure, coordinate frames, and synchronization \cite{rodriguez_mendoza_wearable_2024} \\
    Full-body IMU tracking & Rich body kinematics & Many sensors, calibration, and modeling effort \cite{ruffaldi_sensor_2015} \\
    Low-cost seat/boat IMUs & Simple, cheap, relative movement indicators & Limited to proxy metrics, needs careful interpretation \\
    \hline
  \end{tabular}
\end{table}

The comparison shows why the proposed direction is useful. It is not the most complete measurement system. It is also not the most accurate position system. Its strength is that it is small, cheap, understandable, and focused on relative metrics that fit the limits of low-cost IMUs.

\section{Research Gap}

The literature shows that rowing sensing is already a rich research field. There are systems for boat speed and acceleration, oar angle, handle force, pressure distribution, seat displacement, EMG, crew synchrony, and 3D handle positioning \cite{llosa_remote_2009,avvenuti_mars_2013,cloud_adaptive_2019,hohmuth_wireless_2023,rodriguez_mendoza_wearable_2024}. However, many of these systems are complex, expensive, intrusive, or focused on a different measurement target. Some focus mainly on oars or handles. Others use UWB, GNSS, EMG, multiple IMUs, or contactless distance sensors. These systems can provide rich data, but they increase setup effort and make the system harder to reproduce as a low-cost prototype.

There is also a technical gap between raw IMU signals and reliable rowing metrics. A low-cost IMU can measure acceleration and angular velocity well enough for many short-term patterns, but exact absolute position is hard because of drift and integration errors \cite{cloud_adaptive_2019,rodriguez_mendoza_wearable_2024}. This is especially important for the sliding seat, because a seat-mounted sensor measures both seat motion and boat motion. The two-sensor seat/boat concept by Gravenhorst et al. is therefore a strong starting point \cite{gravenhorst_self-aligning_2012}. The remaining opportunity is to build a simple prototype that uses this relative idea without depending on full double-integrated seat position.

This thesis addresses that opportunity. It develops a low-cost IMU-based prototype for rowing and focuses on relative, interpretable, and stroke-level metrics. The system is not designed to replace laboratory motion capture or high-end rowing instrumentation. Instead, it is designed as a practical step from requirements to prototype evaluation: record useful IMU data, protect the electronics, stream and store the signal, analyze stroke-level patterns, and use a boat reference sensor to reduce the ambiguity of single-sensor seat measurements.

\section{Summary}

The related work shows that rowing measurement systems can be built in many ways. High-end approaches can measure more variables, but they often require more hardware, more calibration, and more setup time. Low-cost IMU approaches are attractive, but they must handle drift, mounting quality, and interpretation limits. Seat motion is especially relevant because it is connected to leg drive and stroke timing, but a single seat sensor mixes seat and boat motion.

The prototype in this thesis is therefore based on a focused design decision: use low-cost IMUs, avoid exact absolute seat position, and derive relative and stroke-level metrics instead. This makes the system less complex than many existing systems while still addressing a meaningful rowing-specific problem.


%%%%%%%%%%%%%%%%%%%%%%%%%%%%%%%%%%%%%%%%%%%%%%%%%%%%%%%%%%%%%%%%%%%%%%%%%%%%%%%
\chapter{Requirements and System Concept} %%%%%%%%%%%%%%%%%%%%%%%%%%%%%%%%%%%%%

\section{Requirement Strategy}

The prototype is developed as a low-cost rowing sensing system. The main goal is not to build a final commercial product, but to build a working research prototype that connects requirements, hardware, firmware, data recording, and analysis. This means that the requirements must be realistic for a bachelor thesis. The system should collect useful data, but it should not claim a level of accuracy that cannot be validated within the project.

The most important requirement is that the system must support rowing-specific motion analysis. Rowing is periodic, dynamic, and sensitive to timing. Therefore, the sensor must record data at a rate that is high enough to describe the stroke cycle. A very low sampling rate would miss peaks and transitions. The prototype therefore uses a target sampling rate of approximately 100 Hz. This is high enough for first stroke-level analysis and still practical for logging to an SD card and streaming over USB or BLE.

The second requirement is low cost. Many related systems use several sensors, GPS, UWB, EMG, force sensors, or commercial rowing equipment. These systems can measure more variables, but they are harder to reproduce. The present prototype should be buildable with common maker hardware and should stay understandable. This makes it easier to adapt, repair, and improve.

The third requirement is unobtrusiveness. The device should not strongly change the rowing situation. It should be small, light, and mounted without long visible cables. This is not only a comfort requirement. It is also a measurement requirement. If the athlete changes posture because the device is visible, heavy, or uncomfortable, the measured movement may no longer represent normal rowing.

The fourth requirement is repeatability. The same movement should produce a similar signal in a later test. For this, the sensor orientation and mounting position must be as stable as possible. The enclosure is therefore part of the sensing system. It must protect the electronics, but it must also keep the IMU fixed relative to the seat or boat.

\section{Functional Requirements}

The prototype must fulfill the following functional requirements:

\begin{itemize}
  \item record acceleration and angular velocity from an IMU;
  \item write the data to a CSV file on an SD card;
  \item include a device identifier, sequence number, and timestamp in each row;
  \item stream the same data over USB serial for debugging and live analysis;
  \item optionally stream data over BLE for wireless use;
  \item support at least one \texttt{SEAT} unit and allow the same firmware to be used for a \texttt{BOAT} reference unit;
  \item provide an offline analysis script for recorded CSV files;
  \item provide a live dashboard for rolling metrics and marker labels;
  \item compute relative signals when seat and boat data are available.
\end{itemize}

The CSV format is intentionally simple. It uses one row per sample and one column per signal. The current format is:

\begin{verbatim}
device_id,sequence,time_us,
acc_x_ms2,acc_y_ms2,acc_z_ms2,
gyro_x_rads,gyro_y_rads,gyro_z_rads
\end{verbatim}

This format is easy to inspect with a spreadsheet, a script, or a text editor. The sequence number helps detect missing samples. The timestamp allows sample timing and sampling frequency to be checked after recording.

\section{Non-Functional Requirements}

The non-functional requirements are especially important because the prototype is meant for a physical sport environment. The system should be:

\begin{itemize}
  \item \textbf{small}, so that it can be mounted without disturbing the athlete;
  \item \textbf{robust}, so that normal handling and vibration do not break the system;
  \item \textbf{splash-resistant}, because rowing takes place near water;
  \item \textbf{easy to start}, because a complicated setup would reduce practical usefulness;
  \item \textbf{transparent}, so that the user can understand what is measured and how metrics are calculated;
  \item \textbf{honest}, so that proxy values are not presented as exact physical quantities.
\end{itemize}

The last point is central for this thesis. A low-cost IMU can produce useful movement signals, but it cannot provide exact position or calibrated boat speed without additional validation. Therefore, the metrics are named and interpreted carefully. Values derived from acceleration integration are called velocity or speed proxies. They describe signal shape, timing, and relative intensity, but they are not claimed to be exact speed in meters per second.

\section{System Concept}

The system is built around two possible sensor roles. The first role is \texttt{SEAT}. This unit is mounted on or near the sliding seat. It records the movement that is closest to the athlete's leg-driven seat motion. However, this signal also contains boat movement because the boat itself accelerates during rowing.

The second role is \texttt{BOAT}. This unit is mounted on the boat shell and acts as a reference. It measures the boat motion without the additional sliding-seat motion. When both units are available, the main relative signal is calculated as:

\begin{equation}
  a_{\text{relative}}(t) = a_{\text{seat}}(t) - a_{\text{boat}}(t).
\end{equation}

This relative signal is not perfect. It depends on sensor alignment, synchronization, and mounting quality. Still, it is a useful design because it reduces the ambiguity of a single seat sensor. Instead of interpreting all seat acceleration as athlete motion, the system can subtract the boat reference and focus on what remains.

The prototype therefore supports two levels of analysis. In one-sensor mode, it can calculate stroke timing, forward acceleration variation, pitch and yaw stability, and a seat velocity proxy. In two-sensor mode, it can additionally calculate SEAT--minus--BOAT acceleration, relative velocity proxies, boat-only velocity proxies, and subtraction effects.

\section{Design Trade-Off}

The chosen design is a compromise. It is less complete than multi-sensor rowing systems with force sensors, GPS, UWB, EMG, or direct seat displacement measurement. It also cannot replace a validated motion capture system. However, it has several advantages. It is cheaper, smaller, easier to build, and easier to understand. It can be tested early and improved step by step.

This trade-off fits the project goal. The thesis is not mainly about creating the most accurate rowing measurement system. It is about developing a practical prototype and showing which useful metrics can be derived from low-cost IMU data while respecting the limits of the sensor.


%%%%%%%%%%%%%%%%%%%%%%%%%%%%%%%%%%%%%%%%%%%%%%%%%%%%%%%%%%%%%%%%%%%%%%%%%%%%%%%
\chapter{Prototype Development} %%%%%%%%%%%%%%%%%%%%%%%%%%%%%%%%%%%%%%%%%%%%%%%

\section{Hardware Platform}

The prototype is based on a small embedded board with an IMU, SD card support, USB serial communication, and BLE communication. The IMU provides three acceleration axes and three gyroscope axes. The acceleration values are used for forward/backward, lateral, and vertical movement analysis. The gyroscope values are used for roll, pitch, and yaw rate indicators.

The hardware choice follows the low-cost requirement. Instead of using a dedicated commercial rowing measurement platform, the prototype uses a general-purpose embedded sensor board. This makes the system flexible. The same firmware can be used for a seat unit or a boat unit by changing the device identifier.

The SD card is important because wireless streaming can fail or become unstable. Local logging gives a more reliable record of the experiment. USB serial is useful during development because it allows direct inspection of the CSV stream. BLE is useful for later live use because it avoids a cable between the rower and the computer.

\section{Firmware}

The firmware samples the IMU at approximately 100 Hz. Each sample receives a sequence number and a timestamp in microseconds. The sample is then written as a CSV row to the SD card and also sent over USB serial. If BLE streaming is enabled and a client is connected, selected rows are also streamed wirelessly.

The sequence number is a simple but useful reliability check. If the sequence numbers are not continuous, samples were lost. This can happen because of SD card delays, communication problems, or software timing issues. The offline analysis therefore checks sequence gaps and reports them.

The timestamp is used to calculate the real sampling interval. The firmware aims for a fixed sample interval, but writing to storage and streaming data can still introduce small timing variations. The analysis script therefore calculates the mean interval and the estimated sampling rate from the recorded file.

\section{Data Pipeline}

The data pipeline has three parts. First, the embedded device produces the raw CSV stream. Second, the live dashboard receives data from USB serial, BLE, or dual BLE. Third, the offline analysis script reads recorded CSV files and generates an HTML report.

The live dashboard is useful during development because it gives immediate feedback. It shows rolling metrics such as sample rate, forward motion, stroke rate, pitch stability, yaw stability, and sequence gaps. In dual-IMU mode, it also shows SEAT, BOAT, and SEAT--minus--BOAT signals. The dashboard can record a segment and allows manual markers such as good or bad strokes.

The offline analysis is useful after recording because it can inspect the full file. It calculates statistics for each signal, detects strokes, estimates stroke rate, creates plots, and calculates relative metrics if a BOAT reference file is available. This is the more stable analysis path because it can process the complete dataset after the experiment.

\section{Enclosure Development}

The enclosure is still an open development task, but it has clear requirements. It must protect the electronics from splashes, water drops, and handling. It must hold the IMU in a stable orientation. It must allow the device to be mounted to the seat or boat without changing position during rowing. It should also allow charging, programming, or SD card access without making the system too fragile.

A first enclosure concept should therefore focus on simple and reliable geometry. The device should sit tightly inside the case so that the board cannot move relative to the housing. If the board moves inside the case, the measured signal may include housing vibration instead of seat or boat motion. Foam or internal spacers can help, but they must not allow the board to wobble.

For sealing, there are several possible options. A fully sealed box is robust but makes USB and SD access harder. A box with a removable lid is easier to use but needs a good seal. A simple prototype can start with a 3D-printed case and a gasket or tape-based splash protection. A later version can use an O-ring, screws, and a more controlled sealing surface.

The mounting concept is as important as the case itself. For the seat unit, the case should be aligned with the seat rail so that the X axis is as close as possible to the forward/backward rowing direction. For the boat unit, the case should be mounted to the shell or a stable boat structure with the same orientation if possible. Clear markings on the enclosure can help repeat the mounting orientation.

The enclosure evaluation should answer practical questions:

\begin{itemize}
  \item Does the electronics board fit securely?
  \item Can the device be switched on, charged, and read out?
  \item Does the case protect against splashes?
  \item Does the mounting stay stable during movement?
  \item Is the sensor orientation repeatable between tests?
  \item Is the case small enough to avoid disturbing the athlete?
\end{itemize}

Even if the final enclosure is not perfect, documenting these questions is valuable. The enclosure is part of the engineering process, and its limitations directly affect signal quality.


%%%%%%%%%%%%%%%%%%%%%%%%%%%%%%%%%%%%%%%%%%%%%%%%%%%%%%%%%%%%%%%%%%%%%%%%%%%%%%%
\chapter{Metrics and Analysis} %%%%%%%%%%%%%%%%%%%%%%%%%%%%%%%%%%%%%%%%%%%%%%%%

\section{Coordinate Interpretation}

The analysis uses a practical mounted-coordinate interpretation. The X acceleration axis is treated as the forward/backward candidate axis, because this direction is most relevant for sliding-seat motion. The Y acceleration axis is treated as lateral movement, and the Z acceleration axis is treated as vertical movement. The gyroscope axes are interpreted as roll, pitch, and yaw rate indicators.

This interpretation depends on mounting. If the sensor is rotated, the axes no longer match the intended rowing directions. Therefore, the enclosure should include orientation markings, and the evaluation should describe how the sensor was mounted.

\section{Basic Signal Metrics}

For every recorded axis, the analysis calculates basic statistics:

\begin{equation}
  \mu = \frac{1}{N}\sum_{i=1}^{N} x_i
\end{equation}

\begin{equation}
  \sigma = \sqrt{\frac{1}{N}\sum_{i=1}^{N}(x_i - \mu)^2}
\end{equation}

\begin{equation}
  \text{RMS} = \sqrt{\frac{1}{N}\sum_{i=1}^{N}x_i^2}
\end{equation}

The mean shows the average signal level. For acceleration, this can include gravity and mounting angle. The standard deviation shows how much the signal varies. The RMS value gives a magnitude indicator when positive and negative values both matter. These values are simple, but they are useful for comparing recordings and checking whether the signal contains strong movement.

\section{Stroke Detection}

Stroke detection is based on the smoothed and centered forward acceleration signal. First, the signal is centered by subtracting a baseline. Then a moving average reduces high-frequency noise. Positive peaks are detected with a minimum time distance between peaks. The number of complete intervals between peaks is used as the stroke count.

The stroke rate is calculated as:

\begin{equation}
  \text{stroke rate} =
  \frac{\text{detected strokes}}{\text{duration in seconds}} \cdot 60.
\end{equation}

This is a first-pass method. It should be validated with video or manual labels in later work. However, it is useful for the prototype because rowing is periodic and produces repeated acceleration patterns.

\section{Velocity and Speed Proxies}

The system estimates velocity proxies by integrating smoothed acceleration:

\begin{equation}
  v_i = v_{i-1} +
  \frac{a_i + a_{i-1}}{2}\Delta t.
\end{equation}

Because IMU integration drifts, the result is corrected by removing a simple linear drift over the analyzed segment. This makes the signal more useful for shape comparison. However, the result is still not a calibrated physical speed. It is therefore called a velocity or speed proxy.

For the boat unit, the boat speed proxy is calculated from the centered and smoothed BOAT forward acceleration. It gives a rough view of the timing and shape of boat acceleration over the recording:

\begin{equation}
  v_{\text{boat,proxy}}(t)
  = \int a_{\text{boat,forward}}(t)\,dt.
\end{equation}

For the two-IMU setup, the relative velocity proxy is calculated from the relative forward acceleration:

\begin{equation}
  v_{\text{relative,proxy}}(t)
  = \int (a_{\text{seat,forward}}(t)
  - a_{\text{boat,forward}}(t))\,dt.
\end{equation}

The analysis also reports when the maximum absolute speed proxy occurs:

\begin{equation}
  \text{peak phase} =
  \frac{t_{\max(|v|)} - t_{\text{start}}}
  {t_{\text{end}} - t_{\text{start}}} \cdot 100.
\end{equation}

This value answers a useful timing question: does the strongest velocity proxy occur early, in the middle, or late in the analyzed window or stroke? In a later rowing evaluation, this can help compare different strokes or different athletes.

\section{Smoothness Metric}

Smoothness is estimated with a jerk-based roughness indicator. Jerk is the rate of change of acceleration:

\begin{equation}
  j_i = \frac{a_i - a_{i-1}}{t_i - t_{i-1}}.
\end{equation}

The report uses the RMS of this jerk signal:

\begin{equation}
  \text{smoothness roughness} =
  \sqrt{\frac{1}{N}\sum_{i=1}^{N}j_i^2}.
\end{equation}

Lower values indicate a smoother acceleration signal. Higher values indicate a more abrupt or noisy signal. This metric is useful because rowing technique is often connected to smooth force application and smooth transitions. However, it is also sensitive to mounting vibration and sensor noise. Therefore, it should be interpreted together with the acceleration plots.

\section{Relative SEAT--BOAT Metrics}

The most important two-sensor metric is the relative forward acceleration:

\begin{equation}
  a_{\text{relative}}(t)
  = a_{\text{seat}}(t) - a_{\text{boat}}(t).
\end{equation}

The analysis also reports a diagnostic sum:

\begin{equation}
  a_{\text{sum}}(t)
  = a_{\text{seat}}(t) + a_{\text{boat}}(t).
\end{equation}

The sum is not the main rowing metric. It is a diagnostic value. If the sum is large, it can indicate shared movement, same-direction motion, or mounting effects. The main signal remains SEAT--minus--BOAT.

The subtraction effect is calculated as:

\begin{equation}
  \text{subtraction effect}
  =
  \left(1 -
  \frac{\sigma_{\text{relative}}}
  {\sigma_{\text{seat}}}
  \right) \cdot 100.
\end{equation}

A positive value means that subtracting the BOAT signal reduced the variation of the raw SEAT signal. A negative value means that the relative signal is stronger than the original SEAT signal. This is not automatically bad. It can happen if the boat and seat move in opposite directions or if the relative seat motion is the dominant movement.

The relative dominance is calculated as:

\begin{equation}
  \text{relative dominance}
  =
  \frac{\sigma_{\text{relative}}}
  {\sigma_{\text{relative}} + \sigma_{\text{boat}}}
  \cdot 100.
\end{equation}

This indicates whether the analyzed window is dominated more by seat-relative movement or by boat reference motion.

\section{Interpretation Limits}

All metrics in this chapter are prototype metrics. They are meant to support interpretation, not to prove final rowing performance. Several limitations remain. The sensors are not hardware-synchronized. The axes depend on mounting. Velocity proxies drift and are not calibrated to true speed. Smoothness can be affected by vibration. Stroke detection can fail if the signal is weak or if the athlete performs non-standard movements.

These limitations do not make the prototype useless. They define the correct interpretation. The system is useful for comparing signal shapes, checking repeatability, detecting strokes, and exploring relative seat-to-boat movement. It is not yet a validated replacement for high-end rowing measurement systems.


%%%%%%%%%%%%%%%%%%%%%%%%%%%%%%%%%%%%%%%%%%%%%%%%%%%%%%%%%%%%%%%%%%%%%%%%%%%%%%%
\chapter{Evaluation Plan} %%%%%%%%%%%%%%%%%%%%%%%%%%%%%%%%%%%%%%%%%%%%%%%%%%%%%

\section{Technical Evaluation}

The first evaluation level is technical. The firmware should compile, start, detect the IMU, create a log file, write a valid CSV header, and record continuous samples. The CSV file should contain increasing sequence numbers and plausible timestamps. The offline analysis should be able to load the file without errors.

The following checks should be reported:

\begin{itemize}
  \item firmware compilation result;
  \item CSV header and example rows;
  \item number of recorded samples;
  \item recording duration;
  \item estimated sampling rate;
  \item number of sequence gaps;
  \item whether USB serial and BLE streaming worked.
\end{itemize}

\section{Signal Evaluation}

The second evaluation level is signal quality. The recorded acceleration and gyroscope signals should be plotted and inspected. A useful recording should show repeated patterns during rowing-like movement. The forward acceleration should contain periodic peaks if strokes are performed. Pitch and yaw rates should show whether the sensor rotates during the motion.

For two-IMU evaluation, the SEAT, BOAT, and relative signals should be compared. The BOAT signal should represent reference motion. The relative signal should show what remains after subtracting the boat motion from the seat motion. The new velocity proxy plots can then show whether the boat-only and relative signals reach their strongest speed proxy at different times.

\section{Enclosure Evaluation}

The enclosure should be evaluated with practical criteria. It should hold the board firmly, protect it from splashes, allow access where needed, and support repeatable mounting. If the enclosure is only a first prototype, this should be stated honestly. The evaluation can still describe what worked, what did not work, and what should be improved.

Important observations include whether the case moved during testing, whether cables or openings were problematic, whether the sensor orientation was clear, and whether the device could be mounted without disturbing the rower.

\section{Limitations of the Evaluation}

The evaluation will likely be limited by time, available equipment, and the maturity of the enclosure. It may not include a full on-water validation with ground truth. This is acceptable if the thesis is clear about the limitation. The goal is to evaluate the prototype as a development step, not to claim final measurement validity.

Future validation should compare the IMU metrics against video, coach labels, commercial rowing instrumentation, or a higher-accuracy reference system. This would make it possible to test whether the detected strokes, speed proxies, smoothness values, and relative metrics correspond to real rowing technique differences.
